%\textit{In this section, I will present and analyse 4 different approached to improve the stability of optimal portfolios. I will keep the description of the first 3 methods short. However, I don't want to remove them completely because I would like to compare the results of those 3 "simpler" methods to the "more complex" Hierarchical Risk Parity method.}

In this section, we present two methods that try to improve the stability of Markowitz's minimum volatility portfolios. The first method, called covariance matrix shrinking, tries to tackle instability by reducing noise in the sample covariance matrix: if the input covariance matrix does not change dramatically across time, then the resulting optimal portfolio won't change dramatically across time either. The second method, called hierarchical risk parity, is an intermediate optimisation method between Markowitz's minimum variance portfolio and the inverse variance portfolio developed by Marcos López de Prado: instead of considering stochastic dependency between the returns of different assets via their covariance directly, assets are classified in clusters according to their correlation first, and then wealth is allocated across the mentioned clusters so that they contribute to the total portfolio risk in equal proportions.

\subsection{Covariance Matrix Shrinking}
    %\textit{This method tries to improve stability by reducing noise in the sample covariance matrix. It is easy to implement and I believe it is also to interpret by analysing the condition number of the shrunk covariance matrix.}
	
	In Markowitz's minimum volatility portfolio construction method, the only input variable that we have is the covariance matrix of asset returns. Hence, the covariance matrix is the only source of instability and errors in the process of portfolio optimisation for this method. We have already shown this in our practical implementation in section \ref{markowitzscursesec}.
	
	In section \ref{markowitzpracticalimplementationsec}, we have used the sample covariance matrix computed using NumPy's \texttt{.cov()} method applied to the matrix containing the time-series of past returns of each of the $N$ investable assets. This method computes the unbiased sample covariance matrix of the past returns' time-series. Hence, the most important parameter for the covariance matrix estimation is the number of past periods returns in our time-series.
	
	As we have already described in section \ref{backtestsettingssec}, we need at least as many past returns' observations as investable assets, so that our covariance matrix can be invertible (even though this does not guarantee invertibility). This is not a problem in our practical implementation, because we consider a relatively small set of investable assets. Nevertheless, one can think of asset managers constructing portfolios using a way bigger set of investable assets, for example equal to the S\&P500 index, that contains almost 500 companies. For them, having enough past returns' observations is not necessarily a simple condition to fulfil for all investable assets. Moreover, we have also shortly discussed the trade-off between long and short look-back periods, i.e. for high and low number of past returns' observations $T$. Table \ref{tab:tvaluetradeoff} summarises the advantages and disadvantages of choosing a low and a high number of past returns' observations to estimate the covariance matrix of asset returns irrespective of the value of $N$.
	
%	\begin{table}[h!]
%	    \centering
%	    \begin{tabular}{c|cc}
%	         & Low $T$ & High $T$ \\
%	         \hline
%	         Advantages & React to market changes quicker/sooner. & More stable optimal portfolio across time, i.e. lower turnover. \\
%	         Disadvantages & blabla & blabla
%	    \end{tabular}
%	    \caption{$T$ Value Trade-off}
%	    \label{tab:tvaluetradeoff}
%	\end{table}
	

    \begin{table}[h!]
        \centering
        \begin{tabularx}{\textwidth}{c|XX}
             & \textbf{Low $T$} & \textbf{High $T$} \\
	         \hline
            Advantages & React to market changes quicker/sooner. & More stable portfolios, lower turnover. \\
            Disadvantages & Less stable portfolios, higher turnover & React to market changes slower/later.
        \end{tabularx}
        \caption{$T$ Value Trade-off}
	    \label{tab:tvaluetradeoff}
    \end{table} 
	
	On the one hand, we can choose a low number of past returns' observations to construct our covariance matrix. Doing this will result in covariance matrices that are more volatile, i.e. do change a lot across time. The estimation will be more sensitive to changes in market regimes and changes in asset returns' relationships, but also to noise, which will lead to more unstable portfolios and to higher portfolio turnover across time. On the other hand, we can also choose a high number of past returns' observations to construct our covariance matrix. This will help to get covariance matrices that are more stable across time, producing optimal portfolios that are also more stable across time and hence reducing turnover. However, this portfolios will also be more inelastic to changes in market regimes and changes in asset returns' relationships. As a practical example, one can think about the changes in performance across the different equity sectors during a business cycle: some sectors like financials tend to perform better during periods of high inflation, while other sectors like IT tend to perform worse \cite{fidelitybusinesscycles}. Investors would like to get optimal portfolios that react to those changes in business cycles in order to profit from the out-performance of some asset classes and sectors, while protecting their portfolios from the under-performance of some other asset classes and sector. Moreover, optimal portfolios constructed on the noise abundant in financial time-series should be avoided, so as to reduce turnover.
	
	Optimising this trade-off is the objective of the covariance matrix shrinking method.
	
	\subsubsection{Ledoit and Wolf}
	    Ledoit and Wolf propose using a convex combination of the sample covariance matrix and a structured estimator, i.e. some matrix containing little noise and which is estimated with high bias. The idea is to reduce the noise contained in the sample covariance matrix by compensating positive and negative error that tends to be contained in its highest entries in absolute value.
    	
    	\begin{definition}[Shrunk Covariance Matrix\cite{ledoit}]\label{shrunkcovariancematrixdefinition}
    	    Let $S \in \mathbb{R}^{N \times N}$ be a sample covariance matrix and $F \in \mathbb{R}^{N \times N}$ some other matrix, also called \textbf{shrinkage target}. The \textbf{shrunk covariance matrix} is given by
    	    \begin{align*}
    	        \Sigma_{shrunk} := \delta F + (1-\delta) S
    	    \end{align*}
    	    where $\delta \in [0, 1 ]$ is called \textbf{shrinkage constant}.
    	\end{definition}
    	
    	In definition \ref{shrunkcovariancematrixdefinition}, we have 2 parameters apart from the sample covariance matrix $\Sigma$ that we need to further specify: the shrinkage target and the shrinkage constant.
    	
    	For the first, the shrinkage target, we need some estimator that can be computed using few parameters, i.e. a structured estimator, and that also contains important information about the measure that we are trying to estimate, i.e. the stochastic dependency between the random variables of the asset returns. In \cite{ledoit2003}, the authors propose using the single-factor matrix of Sharpe as the shrinkage target, which is a single-index asset pricing model \cite{sharpesim}. However, in \cite{ledoit}, the same authors propose a simpler estimator as shrinkage target.
    	
    	\begin{definition}[Constant Correlation Model \cite{ledoit}]
    	    Let $S := [s_{i,j}]_{i,j= \{1, \dots, N \}} \in \mathbb{R}^{N \times N}$ be the sample covariance matrix, then the \textbf{sample constant correlation matrix $F_S$ corresponding to the sample covariance matrix $S$} is defined as
    	    \begin{align*}
    	        F_S := [f_{i,j}]_{i,j= \{1, \dots, N \}} \in \mathbb{R}^{N \times N}, \quad \textbf{where} \quad f_{i,j} := \begin{cases}
    	        s_{i,i} &\quad \text{for } i=j \\
    	        r \sqrt{s_{i,i} s_{j,j}} &\quad \text{for } i \neq j
    	        \end{cases}
    	    \end{align*}
    	    and $r \in \mathbb{R}$ is the average sample correlation, which can be expressed as
    	    \begin{align}\label{averagesamplecorrelation}
    	        r := \frac{2}{(N-1) N} \sum_{i=1}^{N-1} \sum_{j=i+1}^N \frac{s_{i,j}}{\sqrt{s_{i,i} s_{j,j}}}
    	    \end{align}
    	\end{definition}
    	
    	Ledoit and Wolf argue that the constant correlation model gives comparable performance while being easier to implement than the single-factor matrix model of Sharpe \cite{ledoit}. However, they also mention that the constant correlation model is suited for portfolio construction if the set of investable assets contains assets from different classes, for example fixed-income and equities, because it assumes that all pairwise correlations are treated equal for the computation of the average sample correlation $r$. From an economic perspective, this makes sense because the returns of different asset classes can be driven by very different factors: for example, while the stock returns of technology companies are often driven by their margins and customer base growth, the bond returns of the same companies are rather driven by the liquidity of their balance sheet. Nevertheless, we consider this shrinkage target in our practical implementation, because our backtesting period is marked by the market-stress caused by the COVID-19 crisis, and the asset returns during this kind of periods tend to be driven by macroeconomic factors across the different asset classes.
    	
    	The second parameter in definition \ref{shrunkcovariancematrixdefinition}, the shrinkage constant, represents the compromise between the sample covariance matrix and the structured estimator. It can take infinite many different values, even if constrained to be in the interval $[0,1]$ so as to compute the shrunk covariance matrix as a convex combination of the structured estimator and the sample covariance matrix. Ledoit and Wolf propose defining the shrinkage constant as the optimal $\delta^* \in [0,1]$ that minimises the distance between the shrunk covariance matrix and the true covariance matrix, i.e. the covariance matrix that contains the covariances of investable asset returns without noise.
    	
    	\begin{theorem}[\cite{ledoit2003}]\label{optimalshrinkageconstantheorem}
    	        Let $S := [s_{i,j}]_{i,j= \{1, \dots, N \}} \in \mathbb{R}^{N \times N}$ be the sample covariance matrix computed using $T \in \mathbb{N}$ past returns observations, $F_S := [f_{i,j}]_{i,j= \{1, \dots, N \}} \in \mathbb{R}^{N \times N}$ the sample constant correlation matrix corresponding to $S$, $\Sigma := [\sigma_{i,j}]_{i,j= \{1, \dots, N \}} \in \mathbb{R}^{N \times N}$ the true covariance matrix and $F_\Sigma := [\phi_{i,j}]_{i,j= \{1, \dots, N \}} \in \mathbb{R}^{N \times N}$ its constant correlation matrix. Moreover, we define the loss function
    	        \begin{align*}
    	            L(\delta) = \| \delta F_S + (1-\delta) S - \Sigma \|_F
    	        \end{align*}
    	        where $\| \cdot \|_F$ is the Frobenius norm. The shrinkage constant $\delta^*$ that minimises the expected value of $L$ is given by
    	        \begin{align}\label{optimalshrinkageconstant}
    	            \delta^* := \frac{\sum_{i=1}^N \sum_{j=1}^N \mathbb{V}[s_{i,j}] - Cov[f_{i,j}, s_{i,j}]}{\sum_{i=1}^N \sum_{j=1}^N \mathbb{V}[f_{i,j} - s_{i,j}] + (\phi_{i,j} - \sigma_{i,j})^2}
    	        \end{align}
    	        where the function $\mathbb{V}[\cdot]$ designates the variance and $Cov[\cdot, \cdot]$ the covariance.
    	\end{theorem}
    	
    	The proof of this theorem is provided in \cite{ledoit2003} and goes as follows.
    	
    	\begin{proof}
    	        The expected value of the loss function $L$ can be rewritten as
    	        \begin{align*}
    	            \mathbb{E}[L(\delta)] &= \sum_{i=1}^N \sum_{j=1}^N \mathbb{E}[\delta f_{i,j} + (1- \delta) s_{i,j} - \sigma_{i.j}]^2 \\
    	            &= \sum_{i=1}^N \sum_{j=1}^N \mathbb{V}[\delta f_{i,j} + (1 - \delta) s_{i,j}] + (\mathbb{E}[\delta f_{i,j} + (1- \delta) s_{i,j} - \sigma_{i,j}])^2 \\
    	            &= \sum_{i=1}^N \sum_{j=1}^N \delta^2 \mathbb{V}[f_{i,j}] + (1 - \delta)^2 \mathbb{V}[s_{i,j}] + 2 \delta (1 - \delta) Cov[f_{i,j}, s_{i,j}] + \delta^2 (\phi_{i,j} - \sigma_{i,j})^2
    	        \end{align*}
    	        The first and second derivatives of the expected value of the loss function $L$ with respect to $\delta$ are given by
    	        \begin{align*}
    	            \frac{\partial L}{\partial \delta} &= 2 \sum_{i=1}^N \sum_{j=1}^N \delta \mathbb{V}[f_{i,j}] - (1 - \delta) \mathbb{V}[s_{i,j}] + (1- 2 \delta) Cov[f_{i,j}, s_{i,j}] + \delta (\phi_{i,j} - \sigma_{i,j})^2 \\
    	            \frac{\partial^2 L}{\partial \delta^2} &= 2 \sum_{i=1}^N \sum_{j=1}^N \mathbb{V}[f_{i,j} - s_{i,j}] + (\phi_{i,j} - \sigma_{i,j})^2
    	        \end{align*}
    	        Because $ \frac{\partial^2 L}{\partial \delta^2} > 0$ holds, the minimum of the expected value of the loss function $L$ can be computed by setting $\frac{\partial L}{\partial \delta}$ equal to zero. Solving for $\delta^*$ we get
    	        \begin{align*}
    	            \delta^* = \frac{\sum_{i=1}^N \sum_{j=1}^N \mathbb{V}[s_{i,j}] - Cov[f_{i,j}, s_{i,j}]}{\sum_{i=1}^N \sum_{j=1}^N \mathbb{V}[f_{i,j} - s_{i,j}] + (\phi_{i,j} - \sigma_{i,j})^2}
    	        \end{align*}
    	\end{proof}
    	
    	The only problem of the optimal shrinkage constant presented in theorem \ref{optimalshrinkageconstantheorem} is that we don't know the true covariance matrix and can only compute sample variances and covariances. Therefore, Ledoit and Wolf propose using a consistent estimator for $\delta^*$, i.e. an estimator that converges to the true $\delta^*$ for $T \rightarrow \infty$.
    	
    	\begin{theorem}[\cite{ledoit}]
    	    Assuming that asset returns' random variables are independent and identically distributed with finite fourth moment, let $Y := [y_{i,j}]_{i= \{1, \dots, N \}, j= \{1, \dots, T \}} \in \mathbb{R}^{N \times T}$, where $y_{i,j}$ represents the return of the ith asset in the past jth period, and $\bar{y}_i := \sum_{j=1}^T y_{i,j}$. Moreover, we define the matrices $S := [s_{i,j}]_{i,j= \{1, \dots, N \}} \in \mathbb{R}^{N \times N}$ as the sample covariance matrix computed using the returns observations in $Y$, $F_S := [f_{i,j}]_{i,j= \{1, \dots, N \}} \in \mathbb{R}^{N \times N}$ as its sample constant correlation matrix, $\Sigma := [\sigma_{i,j}]_{i,j= \{1, \dots, N \}} \in \mathbb{R}^{N \times N}$ as the true covariance matrix and $F_\Sigma := [\phi_{i,j}]_{i,j= \{1, \dots, N \}} \in \mathbb{R}^{N \times N}$ as its constant correlation matrix. Then
    	    \begin{align*}
    	        \hat{\pi} := \sum_{i=1}^N \sum_{j=1}^N \hat{\pi}_{i,j} \quad \text{where} \quad \hat{\pi}_{i,j} := \frac{1}{T} \sum_{t=1}^T \left( (y_{i,t} - \bar{y}_{i}) (y_{j,t} - \bar{y}_{j}) \right)^2
    	    \end{align*}
    	    is a consistent estimator for $\sum_{i=1}^N \sum_{j=1}^N \mathbb{V}[\sqrt(T) s_{i,j}]$,
    	    \begin{align*}
    	        \hat{\rho} &:= \sum_{i=1}^N \hat{\pi}_{i,i} + \sum_{i=1}^N \sum_{j=1, \; j \neq i}^N \frac{\bar{r}}{2} \left( \sqrt{\frac{s_{j,j}}{s_{i,i}}} \hat{\vartheta}_{i,i,j} + \sqrt{\frac{s_{i,i}}{s_{j,j}}} \hat{\vartheta}_{j,i,j} \right) \\
    	        \text{where} \quad \hat{\vartheta}_{k,i,j} &:= \frac{1}{T} \sum_{t=1}^T \left( (y_{k,t} - \bar{y}_k)^2 - s_{k,k} \right) \left( (y_{i,t} - \bar{y}_i) (y_{j,t} - \bar{y}_j) - s_{i,j} \right)
    	    \end{align*}
    	    and $\bar{r}$ is defined as in equation \ref{averagesamplecorrelation} is a consistent estimator for $\sum_{i=1}^N \sum_{j=1}^N Cov[\sqrt{T} f_{i,j}, \sqrt{T}, s_{i,j}]$ and
    	    \begin{align*}
    	        \hat{\gamma} := \sum_{i=1}^N \sum_{j=1}^N (f_{i,j} - s_{i,j})^2
    	    \end{align*}
    	    is a consistent estimator for $\sum_{i=1}^N \sum_{j=1}^N (\phi_{i,j} - \sigma_{i,j})^2$. Furthermore,
    	    \begin{align*}
    	        \sum_{i=1}^N \sum_{j=1}^N \mathbb{V}[f_{i,j} - s_{i,j}] = O \left( \frac{1}{T} \right)
    	    \end{align*}
    	    Finally, a consistent estimator for the optimal $\delta^*$ defined in equation \ref{optimalshrinkageconstant} is given by
    	    \begin{align}\label{consistentshrinkageconstantestimator}
    	        \hat{\delta}^* := \frac{1}{T} \frac{\hat{\pi} - \hat{\rho}}{\hat{\gamma}}
    	    \end{align}
    	\end{theorem}
    	
    	A proof of this theorem is cannot be found in \cite{ledoit} but is given in \cite{ledoit2003}. The assumption of asset returns' random variables being independent and identically distributed with finite moments is necessary to prove convergence using the Central Limit Theorem. Ledoit and Wolf argue that stock returns do not fulfil this assumption, but that covariance matrix estimators also use this assumption and that introducing extra degrees of freedom in the estimation process to account for dependence and conditional heteroskedastic does not necessarily lead to better out-of-sample performance \cite{ledoit2003}.
    	
    	In \cite{ledoit}, Ledoit and Wolf point to the fact that $\hat{\delta}^* \notin [0,1]$ in equation \ref{consistentshrinkageconstantestimator} may hold. Hence, the final shrinkage intensity proposed by Ledoit and Wolf is defined as
    	\begin{align*}
    	    \delta := \max \left\{ 0, \min \left\{ \frac{1}{T} \frac{\hat{\pi} - \hat{\rho}}{\hat{\gamma}}, 1 \right\} \right\}
    	\end{align*}
    
    \subsubsection{De Nard}
        \todo{decide if this other method should also be described and implemented}
        \cite{denard}
    
    \subsubsection{Practical Implementation}
        \textit{Analyse results comparing the method to Markowitz's method.}
    
%\subsection{Equal Weight and Inverse Variance Allocation}
%    \textit{The 2 methods in this section are naive approaches to portfolio optimisation. Nevertheless, I think they are a good benchmark to compare the rest of methods. Moreover, we can also use them to see if portfolio optimisation is at all useful.}
%    
%    References: \cite{equalweight}, \cite{berndscherer}
%    
%    \subsubsection{Practical Implementation}
%        \textit{Analyse results comparing the method to Markowitz's method.}
    
\subsection{Hierarchical Risk Parity}
    \textit{This subsection, together with subsection \ref{nassetssec}, is the biggest part of the thesis. I will describe the method in a more mathematical way, as we already discussed. Moreover, I will discuss the differences between this method and the original method of Markowitz, the model assumptions and also its limitations.}
    
    References: \cite{lopezmain}, \cite{berndscherer} \cite{stefanjansen}
    
    \subsubsection{Practical Implementation}
        \textit{Analyse results comparing the method to Markowitz's method and also to the other 3 methods above.}